\section*{Part VI: Human Learning in an Age of Abundant Information and AI}

\subsection*{79. When Information Becomes Cheap, Selection and Judgment Become Scarce}

A learner can access textbooks, online courses, documentation, lectures, and AI-generated explanations at little or no marginal cost. This changes the commercial role of an educational institution. Mere access to a definition or a sequence of topics is less distinctive than it once was.

But abundant information does not imply abundant competence. Learners still need to know what to study, interpret explanations through prerequisites, choose practice, recognize errors, sustain effort, and transfer knowledge to real tasks. The economic opportunity may shift toward \textbf{diagnosis, disciplined practice, feedback, trustworthy assessment, and a social environment that sustains difficult work}.

A new institute should not assume it can sell generic content at an enormous premium merely because the website looks academic. Its stronger potential advantage lies in organizing a transformation that independent learners find difficult to reproduce alone.

\subsection*{80. AI Creates a Distinction Between Doing the Work and Becoming Able to Do the Work}

Imagine a student asks an AI model to produce a functioning application. The model writes most of the code, supplies tests, and explains the architecture. The student presents a polished demo. Has learning occurred?

Possibly. The learner may have compared alternatives, identified bad outputs, adjusted requirements, and learned system judgment. Or they may have copied an artifact they cannot explain. The final product alone cannot distinguish these cases.

A 2025 field experiment in high-school mathematics compared differently designed AI assistance. Access improved performance while tools were available, but relatively unguarded assistance could reduce later unaided performance; learning-oriented safeguards mitigated the harm in that context [R16]. This is not a universal verdict against AI. It demonstrates why \textbf{assisted task performance and acquired independent capability are different outcomes}.

\subsection*{81. AI Tutors Are Designs, Not a Single Category of Intervention}

An AI tutor can provide direct answers, graded hints, questions, feedback, simulations, examples, practice schedules, and conversational support. These are different pedagogies. The prompt, task design, feedback policy, subject, learner, and assessment condition influence what happens.

A randomized study in a college physics setting found stronger immediate performance under a specially designed AI-tutoring experience than under the particular in-class active-learning comparison studied [R17]. The result is promising, but it does not establish that generic AI chats outperform all teachers, all classrooms, or longer-term learning across disciplines.

The founder should ask: \emph{What cognitive work does the system induce or replace?} A chatbot that answers instantly may be convenient. A tutor that helps the learner make the next attempt can serve a different, sometimes more valuable purpose.

\subsection*{82. The Fading-Scaffold Principle for AI}

A plausible learning-oriented AI interaction has 6 phases: \textbf{attempt 1st $\to$ identify where the learner is stuck $\to$ provide the least sufficient hint $\to$ request a revised attempt $\to$ test with a new variant $\to$ fade the help}.

For example, when a student asks how to fix leakage in an evaluation pipeline, the tutor might 1st ask which fields are available at prediction time. If the learner cannot answer, it can show a timeline and a narrower question. Only after continued difficulty should it reveal a fuller worked example. The next task should require the student to recognize the same principle without the scaffolding.

This is not a prohibition on using AI for production work. It is a distinction between \textbf{training a skill} and \textbf{getting an immediate artifact}. The right design depends on which of those is the educational objective.

\subsection*{83. Verification Is the New Basic Literacy}

When AI can produce fluent, plausible answers, learners must become skilled at deciding whether the answer should be trusted. In engineering, this involves tests, source inspection, permissions, numerical checks, evaluation sets, and understanding failure modes. In research, it involves provenance, citation inspection, alternative hypotheses, and methods that make claims falsifiable.

An AI-intensive curriculum should require students to generate \textbf{verification plans}, not just prompts. Before accepting a generated solution, they should be able to say which failure would matter, what test could reveal it, and what evidence would justify use.

This is an opportunity for a research-active educational institution: to cultivate people who can work productively with powerful tools without surrendering responsibility for truth and consequences.

\subsection*{84. The Correct Human-AI Boundary Depends on the Competence Claim}

For a student learning algebraic fluency, AI completing every step removes core practice. For an advanced analyst working on a real modeling project, using AI to write boilerplate code may free attention for experimental design. For a surgeon-in-training, an AI explanation cannot substitute for the physical performance being assessed.

Define the competency, then classify tool use as \textbf{essential}, \textbf{permitted}, \textbf{limited}, or \textbf{prohibited for that assessment}. Make the policy educationally intelligible. An arbitrary ban can misrepresent authentic professional work; unlimited hidden assistance can make assessment meaningless.

A mature program often uses both unaided probes of core knowledge and authentic assisted tasks testing tool selection, verification, and judgment. It should not confuse 1 with the other.

\subsection*{85. The Risk of AI-Induced Dependence}

If learners can obtain relief from discomfort by asking for a completed answer, they may rationally take the shortcut. This is not necessarily a character defect. The learning environment may be rewarding completion, speed, and grades more strongly than independent reasoning.

Redesign incentives: grade consequential reasoning, include unseen tasks, allow revision without immediate penalty, and use graded hints. Avoid building a hostile surveillance culture in which every legitimate use of assistance is treated as cheating. Instead, make authorship and assistance expectations explicit and assess the skills the program claims to teach.

AI can multiply opportunities for practice if designed well. It can also multiply convincing evidence of skills the learner does not possess. The founder should treat both as real possibilities.
